\documentclass[conference]{IEEEtran}
\usepackage{cite}
\usepackage{amsmath,amssymb,amsfonts}
\usepackage{graphicx}
\usepackage{textcomp}
\usepackage{xcolor}
\usepackage{booktabs}
\usepackage{url}
\usepackage[hidelinks]{hyperref}
\usepackage{tagging}
\usepackage{sc26repro}

\begin{document}
\twocolumn[%
{\begin{center}
\Huge Appendix: Artifact Description/Artifact Evaluation
\end{center}}
]

\appendixAD
\section{Overview of Contributions and Artifacts}

\subsection{Paper's Main Contributions}
\begin{description}
\item[$C_1$] Reactive Subflow Spraying (RSS), an end-host load balancer that
routes a small number of ordered subflows and periodically moves the
worst-performing subflow using end-host congestion signals.
\item[$C_2$] Precise Fast Loss Detection (PFLD), a receiver-assisted recovery
mechanism based on per-subflow sequence numbers, selective negative
acknowledgments, and proactive and tail probes.
\item[$C_3$] A low-state RSS and PFLD co-design that combines multipath load
balancing with precise recovery and requires no switch modification.
\item[$C_4$] Evaluation using packet-level simulation, microbenchmarks, AI
collectives, LLaMA and HPC execution traces, sensitivity experiments, and a
programmable-DPU prototype.
\end{description}

\subsection{Computational Artifacts}
The evaluator repository is
\url{https://github.com/tommasobo/SC26_OrderedChaos}. The archived artifact is
available at \url{https://doi.org/10.5281/zenodo.21542397}. The repository
provides two top-level reproduction commands. By default, generated working
data stays inside the clone and final plots are collected under
\path{results/}. An optional Slurm backend can place the same data on
high-capacity storage.

\begin{center}
\footnotesize
\begin{tabular}{@{}p{0.13\columnwidth}p{0.20\columnwidth}p{0.57\columnwidth}@{}}
\toprule
Artifact & Contrib. & Paper elements \\
\midrule
$A_1$ & $C_1$ to $C_4$ & htsim source, tests, experiment manifest, public-input
fetcher, local and optional Slurm workflows, analysis scripts, Figures 1 and
5 to 13, Table I,
Camera Ready New Image 1, and Camera Ready New Image 2. \\
\midrule
$A_2$ & $C_2$, $C_4$ & Figure 14 plotting inputs and scripts for the DPU
summary measurements. \\
\bottomrule
\end{tabular}
\end{center}

Figures 2 to 4 are protocol and design illustrations and therefore require no
computational workflow.

\section{Artifact Identification}
\newartifact
\artrel
Artifact $A_1$ contains the simulator implementation of RSS, PFLD, RACK-TLP,
and the routing baselines. The machine-readable manifest fixes the complete
configuration for every result. Each runner requires a new output directory
and records the literal command, executable hash, input hash, runtime, memory,
and pass status. Packaged tables from the paper are not used as experimental
input.

\artexp
The expected result is a new plot or table with the same axes and algorithms as
the corresponding paper element. In particular, RSS should remain competitive
with the multipath baselines under congestion, and RSS+PFLD should reduce loss
recovery delay and exposed communication time when packets are lost.

The two supplementary camera-ready results use names that remain separate from
the submitted paper numbering. Camera Ready New Image 1 presents Google RPC
FCT distributions as violin plots. Camera Ready New Image 2 presents Google
RPC timeout-event counts.

Completed runs collect numbered paper PDFs and PNGs in the result directory
and all supplementary camera-ready PDFs and PNGs in its
\path{ExtraFigures/} subdirectory. The result directory also contains a
\texttt{PASS} marker and the execution identifier used to generate it.

\arttime
For local planning, we assume a recent laptop with 8 CPU cores, 32 GiB of
memory, a local SSD, and the default two workers. Under this assumption, allow
15 to 30 minutes for the short workflow, including environment provisioning
and a clean build, and 24 to 48 hours for the full workflow. These laptop
values are estimates and depend on processor speed and available memory.

The table below gives measured job wall times and resource requests for the
optional single-node Slurm backend. Independent Slurm jobs may run concurrently,
so its end-to-end time is the longest dependency chain rather than the sum of
all rows. Allow about 10 minutes for setup and about 90 minutes for the full
critical path, excluding queue delay and the optional LaTeX build. These values
are rounded upward from the observed runs. No experiment uses a GPU.
These are parallel-run timings, not laptop estimates. The default local runner
uses two workers and executes the large experiment groups sequentially, so its
full workflow is substantially slower.

\begin{table*}[t]
\centering
\footnotesize
\begin{tabular}{@{}lrrrrl@{}}
\toprule
Result & CPU & Memory & Full wall time & Scratch & Quick mode \\
\midrule
Build, tests, smoke & 16 & 64 GiB & 5 min & $<0.1$ GiB & same \\
Public traces & 2 & 4 GiB & 5 min & 1.8 GiB & omit \\
Figure 1A & 8 & 32 GiB & 5 min & 2.5 GiB & omit \\
Figure 1B & 24 & 128 GiB & $<1$ min & $<0.1$ GiB & reduced \\
Figure 5, Table I & 1 & 4 GiB & 1 min & $<0.1$ GiB & same \\
Figures 6A/6B & 128 & 64 GiB & 15 min & 0.4 GiB & reduced \\
Figures 7A/7B & 128 & 128 GiB & 75 min shared & 0.4 GiB & reduced \\
Figure 8 & 16 & 96 GiB & 30 min & 4.7 GiB & omit \\
Figure 9 & 8 & 64 GiB & 5 min & $<0.1$ GiB & omit \\
Figure 10 & 32 & 128 GiB & 60 min & 0.3 GiB & reduced \\
Figure 11 & 48 & 192 GiB & 5 min & 2.7 GiB & omit \\
Figure 12 & 48 & 128 GiB & 30 min shared & 21 GiB & omit \\
Figures 13A/13B & 64 & 128 GiB & 15 min & 0.1 GiB & reduced \\
Figure 14 & 1 & 4 GiB & $<1$ min & $<0.1$ GiB & omit \\
Camera Ready New Images 1/2 & 12 & 128 GiB & 15 min shared & 3.1 GiB shared & omit \\
\bottomrule
\end{tabular}
\caption{Evaluator resource guide for the optional Slurm backend. CPU is
requested \texttt{cpus-per-task}.
All jobs record wall time, peak memory, literal commands, and output
provenance in the selected work directory. Times are conservative planning
values rounded up from observed parallel runs. They should not be used to
estimate the sequential two-worker laptop workflow. Under the laptop
assumption above, budget 15 to 30 minutes for the short run and 24 to 48 hours
for the full run. Simulator performance depends on processor model and
effective CPU allocation, so Figure 10 and the micro-collective jobs can vary
substantially across systems.}
\label{tab:runtime}
\end{table*}

The reviewer quick subset focuses on Figure 1B, Figures 6A/6B, Figures 7A/7B,
Figure 10, Figure 5, and Table I. These compact result families exercise the
simulator, loss-recovery mechanisms, plotting, and analytical model. The
optional parallel workflow should be allocated about 60 minutes, excluding
scheduler queue time.

\artin
\artinpart{Hardware}
A Linux laptop or workstation is sufficient for the default workflow. We
recommend 16 GiB of memory for the short path and 32 GiB plus 100 GiB of free
storage for the sequential full path. The optional Slurm backend uses one
CPU-only node per task; requests range from 4 to 192 GiB according to the
experiment. No multi-node allocation, GPU, or special interconnect is required
by the simulator.

\artinpart{Software}
The workflow supports Linux on x86-64 and 64-bit Arm. The current end-to-end
test used 64-bit Arm. The required software and tested versions are:
\begin{itemize}
\item OrderedChaos htsim, GCC/G++ 12.3.0, and CMake 3.28.3. htsim requires
C++17 and CMake 3.16 or newer.
\item Python 3.13.5 and uv 0.9.24. The required package pins are NumPy 2.3.4,
pandas 2.3.3, Matplotlib 3.10.8, Seaborn 0.13.2, SciPy 1.16.3, and PyMuPDF
1.27.1.
\item Git 2.43.0, GNU Bash 4.4.23, GNU time 1.9, GNU coreutils 8.32, and curl
8.6.0.
\item Optional parallel execution uses Slurm 25.05.4. No scheduler is needed
for the default workflow.
\end{itemize}
Exact Python pins are in \path{artifact/requirements.txt}. Tectonic 0.15.0
is needed only to rebuild this appendix; the compiled PDF is included.

\artinpart{Datasets / Inputs}
Synthetic incast, tornado, ring, and all-to-all matrices and all network
topologies are versioned in the repository. LLaMA, ICON, LAMMPS, and MILC
binary traces are fetched from the public SPCL trace repository at
\url{http://storage2.spcl.ethz.ch/traces/}. The fetch job records and verifies
every SHA-256 checksum. The measured download size is 1.8 GiB. If the public
server is temporarily unavailable, the workflow records the failed input,
skips only its dependent trace figure, and continues all independent
experiments. The result directory then contains
\texttt{Figure\_11\_SKIPPED.txt} or
\texttt{Figure\_12\_SKIPPED.txt} with the reason. Re-running the full command
when the server is available produces the omitted figure. Figure 14 uses three
small, versioned CSV files containing the recorded DPU summary values.

\artinpart{Installation and Deployment}
Clone the repository and run \path{reproduce.sh} in either \texttt{short} or
\texttt{full} mode. The default local helper installs the specified uv, Python,
and package versions under \path{.orderedchaos-work/}, builds the simulator,
runs all 43 tests and the smoke simulation, and then runs each required
experiment with laptop-safe parallelism. The optional \texttt{--slurm} backend
accepts account and partition values at launch time through environment
variables; no site-specific value is stored in the repository.

\artcomp
The computational dependency graph is
$T_1\rightarrow(T_3,T_4)\rightarrow T_5$ and
$T_2\rightarrow T_4$, where:
\begin{description}
\item[$T_1$] Cleanly compile htsim and run all tests and the smoke case.
\item[$T_2$] Fetch and checksum public LLaMA and HPC inputs. This can overlap
with $T_1$.
\item[$T_3$] Regenerate Figure 5, Table I, and Figure 14 from versioned inputs.
\item[$T_4$] Run independent fresh simulator matrices. Every large matrix is
split by result and may execute concurrently after $T_1$ and, for trace jobs,
$T_2$.
\item[$T_5$] Parse raw logs, calculate medians and percentiles, create plots,
and build the visual result index.
\end{description}

The full and short configurations, including repetitions and random seeds, are
defined in the manifest. Figure 1B pools matched loss events and uses a
seed-stratified bootstrap interval. Error bars elsewhere are generated from
the repeated runs and are never synthesized.

\artout
Raw simulator logs, command files, timing records, and metrics are written to
\path{.orderedchaos-work/orderedchaos_ae} by default. The final PDFs, PNGs, and
tables are collected under \path{results/}. A completed matrix has a
\texttt{run\_status.csv} whose required rows all read \texttt{pass}. Plotting
commands consume the fresh job-specific root recorded in provenance.

\newartifact
\artrel
Artifact $A_2$ reconstructs the Figure 14 presentation from explicit DPU
summary measurements and supports the probe-overhead and probe-FCT portions of
$C_2$ and $C_4$.

\artexp
The expected summary reports a maximum probe overhead of 6.352\%, a 1-MiB
all-probe speedup of 2.483$\times$, retransmissions equal to corruption drops,
and zero all-probe timeouts. The workflow validates the CSV schemas and creates
both PDF and PNG outputs.

\arttime
Setup is shared with $A_1$ and adds no extra time. Schema validation, metric
calculation, plot rendering, and JSON analysis together take less than
1 minute. The complete $A_2$ workflow uses one CPU and less than 128 MiB of
memory.

\artin
\artinpart{Hardware}
No special hardware is needed to evaluate the supplied summary and plotting
workflow.
\artinpart{Software}
The same pinned Python environment as $A_1$ is used.
\artinpart{Datasets / Inputs}
Inputs are \path{artifact/data/figure14_probe_overhead.csv},
\path{artifact/data/figure14_probe_fct.csv}, and
\path{artifact/data/figure14_probe_correctness.csv}.
\artinpart{Installation and Deployment}
No deployment step beyond the Python environment is required.

\artcomp
The script validates the three inputs, computes the headline quantities, and
renders both Figure 14 panels.

\artout
The output directory contains PDF, PNG, and JSON files. The JSON provides the
numeric checks used by the evaluator.

\appendixAE
\arteval{1}
\artin
The commands below run directly on a Linux laptop or workstation. They require
no scheduler, account name, GPU, or external storage setting.

\begin{flushleft}
\ttfamily
git clone \textbackslash\\
\url{https://github.com/tommasobo/SC26_OrderedChaos}\\
cd SC26\_OrderedChaos
\end{flushleft}

\artcomp
For the fastest evaluator path, run:

\begin{verbatim}
./reproduce.sh short
\end{verbatim}

This runs fresh reduced versions of Figure 1B, Figure 6, Figures 7A/7B, and
Figure 10, and regenerates Figure 5 and Table I. The command prints the final
result directory. Inspect its \texttt{PASS} marker and clearly named plots.
The quick path checks functionality and the main trends. The full workflow is
used for final numerical assessment.

Run the complete workflow with:

\begin{verbatim}
./reproduce.sh full
\end{verbatim}

The local wrapper executes the complete workflow in dependency order. Every
stage creates a unique run-specific root and refuses cached experiment output.
The machine-readable manifest records the complete scientific configuration.
Set \texttt{ORDEREDCHAOS\_JOBS} to adjust local worker concurrency; this
changes wall time but not the scientific configuration.

For higher concurrency, a Slurm installation can execute independent stages
concurrently and is recommended for the complete workflow:

\begin{verbatim}
export SCRATCH=/path/to/high-capacity/storage
./reproduce.sh short --slurm
./reproduce.sh full --slurm
\end{verbatim}

If required by the site, supply account and partition arguments at launch
through the optional environment variables documented in the repository
README.

\artout
For each result, verify the following:
\begin{enumerate}
\item Every required row in \texttt{run\_status.csv} is \texttt{pass}.
\item The recorded seed count and task count agree with the manifest.
\item The plot provenance names the new job-specific raw root.
\item The PDF or PNG contains the expected algorithms, trim modes, and axes.
\item The generated table reports the expected qualitative trend and its
measured median or percentile interval.
\end{enumerate}

Raw output is under \path{.orderedchaos-work/orderedchaos_ae/raw}; intermediate
summaries are under the adjacent \path{analysis} directories; and the clearly
named paper plots are under \path{results/}. If a public binary trace is
unavailable, inspect \texttt{SKIPPED\_RESULTS.txt} and the corresponding
per-figure skip record; all other required status rows and plots remain
independently verifiable. To clean a campaign, preserve the desired final
result directory and remove \path{.orderedchaos-work/}. Optional Slurm runs
place the same structure below \texttt{\$SCRATCH}.

\arteval{2}
\artin
Activate the same Python environment used for $A_1$. No DPU allocation is
required for the supplied Figure 14 summary workflow.

\artcomp
Run the plotter with a new output directory:

\begin{verbatim}
.orderedchaos-work/orderedchaos_ae/venv/bin/python \
  artifact/scripts/plot_figure14_summary.py \
  --output results/figure14-check
\end{verbatim}

\artout
Expected output is PDF and PNG plus JSON reporting 6.352\% maximum overhead,
2.483$\times$ speedup at 1 MiB, retransmissions equal to drops, and zero
all-probe timeouts.

\end{document}
